%========================================
% MatinBook — Stage 12: Bibliography Test
% Version: v1.1
% Purpose:
%   Validate the bibliography system of MatinBook v1.1:
%     1. \addbibresource with an external .bib file
%     2. \cite, \parencite, \textcite, \autocite
%     3. Multiple citations in one command
%     4. \printbibliography with a Persian title
%     5. Full biber compilation cycle
%     6. Persian bibliography entries
%
% Compilation:
%   xelatex -shell-escape stage12-biblatex.tex
%   biber stage12-biblatex
%   xelatex -shell-escape stage12-biblatex.tex
%   xelatex -shell-escape stage12-biblatex.tex
%   (biber is required to generate the bibliography)
%
% Expected outcome:
%   A professionally typeset PDF demonstrating that all
%   citation commands work correctly, that the bibliography
%   is generated with the correct Persian title and RTL
%   layout, and that no "??" markers appear.
%========================================

%------------------------------------------------------------------------
% Search Path (for tests directory)
%------------------------------------------------------------------------
\makeatletter
\def\input@path{{../}}
\makeatother

%------------------------------------------------------------------------
% Document Class
%------------------------------------------------------------------------
\documentclass{matinbook}

%------------------------------------------------------------------------
% Bibliography Resource
%------------------------------------------------------------------------
\addbibresource{../references.bib}

%------------------------------------------------------------------------
% Book Metadata
%------------------------------------------------------------------------
\title{آزمون سیستم منابع و کتاب‌نامه}
\author{تیم توسعه MatinBook}
\date{مهر ۱۴۰۵}

\booktitle{آزمون سیستم منابع و کتاب‌نامه}

\renewcommand{\repository}{github.com/matinmahmoudi-cs/matinbook}

%========================================================================
% DOCUMENT
%========================================================================
\begin{document}

%------------------------------------------------------------------------
% FRONT COVER
%------------------------------------------------------------------------
\makecover
    {آزمون سیستم منابع و کتاب‌نامه}
    {ارزیابی استناد و کتاب‌شناسی در MatinBook}
    {تیم توسعه MatinBook}
    {مهر ۱۴۰۵}

%------------------------------------------------------------------------
% FRONT MATTER (symmetric margins)
%------------------------------------------------------------------------
\frontmatter
\frontmattergeometry

% Title page
\maketitle

% Copyright / Colophon
\thispagestyle{empty}
\vfill
\begin{center}
    {\large\textbf{شناسنامه آزمون}}

    \vspace{0.8cm}

    \begin{tabular}{rl}
        عنوان: & آزمون سیستم منابع و کتاب‌نامه \\
        نسخه: & \matinbookversion \\
        تاریخ: & \matinbookdate \\
        موتور: & \lr{XeLaTeX} \\
        مجوز: & \lr{MIT} \\
    \end{tabular}

    \vspace{0.8cm}

    {\small این سند بخشی از مجموعه آزمون‌های یکپارچگی \lr{MatinBook} است.}
\end{center}
\clearpage

% Preface
\chapter*{پیشگفتار}
\addcontentsline{toc}{chapter}{پیشگفتار}

این سند، دوازدهمین آزمون از مجموعه‌ی آزمون‌های یکپارچگی
\lr{MatinBook} نسخه‌ی \lr{\matinbookversion} است. هدف آن،
اعتبارسنجی کامل سیستم کتاب‌نامه‌ی کلاس در شش حوزه‌ی
متمایز است: بارگذاری فایل \lr{.bib} خارجی، دستورات
استناد (\lr{\textbackslash cite}، \lr{\textbackslash parencite}،
\lr{\textbackslash textcite}، \lr{\textbackslash autocite})،
استنادهای چندگانه، نمایش کتاب‌نامه با
\lr{\textbackslash printbibliography}، چرخه‌ی کامل
\lr{biber}، و ورودی‌های کتاب‌نامه‌ی فارسی.

در کتاب‌های علمی، مدیریت منابع و کتاب‌نامه از ارکان
اساسی نگارش است. هر ادعای علمی باید به منبع معتبر
ارجاع داده شود و خواننده باید بتواند به‌راحتی منابع
را پیدا کند. \lr{MatinBook} با استفاده از \lr{biblatex}
و موتور \lr{biber}، این قابلیت‌ها را فراهم می‌سازد.

هر بخش از این آزمون، با مثال‌های عملی و ارجاعات متقابل
همراه است تا خواننده بتواند درستی رفتار هر زیرسیستم را
در بافت واقعی یک کتاب ارزیابی کند.

\clearpage

% Table of Contents
\tableofcontents
\clearpage

% List of Figures
\listoffigures
\clearpage

% List of Tables
\listoftables
\clearpage

%------------------------------------------------------------------------
% MAIN MATTER (wide margin for notes)
%------------------------------------------------------------------------
\mainmatter
\mainmattergeometry

%========================================================================
% CHAPTER 1 — INTRODUCTION TO CITATIONS
%========================================================================
\chapter{مقدمه‌ای بر استناددهی}
\label{chap:citation-intro}

\section{اهمیت استناددهی}
\label{sec:citation-importance}

استناددهی\index{استناددهی} یکی از ارکان اساسی نگارش
علمی است. هر کتاب علمی باید منابع مورد استفاده‌ی خود
را به‌طور دقیق مشخص کند. \lr{MatinBook} از بسته‌ی
\lr{biblatex} با موتور \lr{biber} برای مدیریت منابع
استفاده می‌کند.

\mnote{استناددهی دقیق، اعتبار علمی کتاب را افزایش می‌دهد و از سرقت ادبی جلوگیری می‌کند.}

\section{شیوه‌ی استناد}
\label{sec:citation-style}

در این کتاب از شیوه‌ی \textbf{عددی} (\lr{numeric}) برای
استناددهی استفاده می‌شود. هر منبع با یک شماره مشخص
می‌شود و در بخش کتاب‌نامه به ترتیب استناد فهرست
می‌شود.

%========================================================================
% CHAPTER 2 — CITATION EXAMPLES
%========================================================================
\chapter{نمونه‌های استناد}
\label{chap:citation-examples}

\section{کتاب‌ها}
\label{sec:citation-books}

کتاب‌های مرجع در زمینه‌ی علوم کامپیوتر:

\begin{itemize}
    \item \textbf{هنر برنامه‌نویسی:} دونالد کنوث
    \cite{knuth1984} مجموعه‌ای بی‌نظیر درباره‌ی
    الگوریتم‌ها و ساختمان داده‌ها نوشته است.

    \item \textbf{راهنمای لاتک:} لسلی لمپورت
    \cite{lamport1994} سیستم \LaTeX{} را برای
    حروف‌چینی متون علمی معرفی کرد.

    \item \textbf{مقدمه‌ای بر الگوریتم‌ها:} کتاب
    \lr{CLRS} \cite{cormen2009} مرجع استاندارد
    الگوریتم‌ها در دانشگاه‌های سراسر جهان است.
\end{itemize}

\section{مقالات علمی}
\label{sec:citation-articles}

\begin{itemize}
    \item \textbf{نسبیت خاص:} آلبرت اینشتین در سال
    ۱۹۰۵ مقاله‌ی معروف خود درباره‌ی الکترودینامیک
    اجسام متحرک را منتشر کرد \cite{einstein1905}.
    این مقاله پایه‌گذار نظریه‌ی نسبیت خاص بود.
\end{itemize}

\section{منابع آنلاین}
\label{sec:citation-online}

\begin{itemize}
    \item \textbf{پروژه لاتک:} وبسایت رسمی پروژه
    \LaTeX{} \cite{latex-project} منبع معتبری برای
    دریافت اطلاعات به‌روز است.
\end{itemize}

\section{استنادهای چندگانه}
\label{sec:citation-multiple}

مطالعه‌ی هم‌زمان چند منبع برای درک عمیق یک موضوع
ضروری است. برای مثال، برای یادگیری الگوریتم‌ها
می‌توان به آثار کنوث \cite{knuth1984} و کورمن
\cite{cormen2009} مراجعه کرد.

همچنین برای آشنایی با \LaTeX{}، منابع
\cite{lamport1994,latex-project} توصیه می‌شوند.

%========================================================================
% CHAPTER 3 — CITATION COMMANDS
%========================================================================
\chapter{دستورات استناد}
\label{chap:citation-commands}

\section{دستورات پایه}
\label{sec:citation-base}

\lr{MatinBook} از تمام دستورات استاندارد \lr{biblatex}
پشتیبانی می‌کند. جدول \cref{tab:citation-commands}
این دستورات را نشان می‌دهد:

\begin{matintable}{دستورات استناد در \lr{biblatex}}{tab:citation-commands}
\begin{tblr}{
    width = \textwidth,
    colspec = {l X[c] X[c]},
    hlines, vlines,
    colsep = 8pt,
    rowsep = 4pt,
    row{1} = {font=\bfseries},
}
دستور & توضیحات & مثال \\
\lr{\textbackslash cite} & استناد ساده & \cite{knuth1984} \\
\lr{\textbackslash parencite} & استناد پرانتزی & \parencite{lamport1994} \\
\lr{\textbackslash textcite} & استناد متنی & \textcite{cormen2009} \\
\lr{\textbackslash autocite} & استناد خودکار & \autocite{einstein1905} \\
\end{tblr}
\end{matintable}

\section{تنظیمات کتاب‌نامه}
\label{sec:citation-settings}

تنظیمات \lr{biblatex} در \lr{MatinBook}:

\begin{itemize}
    \item \textbf{موتور:} \lr{biber}
    \item \textbf{شیوه:} \lr{numeric} (عددی)
    \item \textbf{ترتیب:} \lr{none} (به ترتیب استناد)
    \item \textbf{فایل منبع:} \lr{references.bib}
\end{itemize}

\section{ارجاع به منابع بدون استناد}
\label{sec:citation-nocite}

می‌توان منابعی را که مستقیماً به آن‌ها استناد نشده
اما مطالعه شده‌اند، با \lr{\textbackslash nocite}
در کتاب‌نامه گنجاند:

\begin{latin}
\begin{minted}{latex}
\nocite{*}
\end{minted}
\end{latin}

> **توجه:** این دستور در این آزمون فعال نیست، چون
> می‌خواهیم فقط منابع استنادشده نمایش داده شوند.

%========================================================================
% CHAPTER 4 — BIBLIOGRAPHY DISPLAY
%========================================================================
\chapter{نمایش کتاب‌نامه}
\label{chap:bibliography-display}

\section{فرمت خروجی}
\label{sec:bibliography-format}

کتاب‌نامه در \lr{MatinBook} با عنوان «منابع و مراجع»
و به‌صورت راست‌به‌چپ نمایش داده می‌شود. برخلاف نسخه‌های
قدیمی، \lr{\textbackslash printbibliography} **نباید**
داخل محیط \lr{latin} قرار گیرد، زیرا \lr{xepersian}
خودش جهت متن را به‌درستی مدیریت می‌کند.

\mnote{در v1.1، \lr{\textbackslash printbibliography} را بدون \lr{latin} بنویسید.}

\section{چرخه‌ی کامپایل}
\label{sec:bibliography-compile}

برای به‌روزرسانی کتاب‌نامه پس از افزودن استناد جدید،
باید دستور \lr{biber} را اجرا کنید:

\begin{latin}
\begin{minted}{bash}
xelatex -shell-escape document.tex
biber document
xelatex -shell-escape document.tex
xelatex -shell-escape document.tex
\end{minted}
\end{latin}

\begin{remark}
    اگر کتاب‌نامه خالی است یا استنادها به‌صورت «؟؟»
    نمایش داده می‌شوند، یعنی \lr{biber} اجرا نشده است.
    چرخه‌ی کامل بالا را دوباره اجرا کنید.
\end{remark}

%========================================================================
% CHAPTER 5 — SUMMARY
%========================================================================
\chapter{خلاصه‌ی نتایج}
\label{chap:summary}

\section{وضعیت سیستم کتاب‌نامه}
\label{sec:summary-status}

جدول \cref{tab:bibliography-status} وضعیت قابلیت‌های
کتاب‌نامه‌ی \lr{MatinBook} را نشان می‌دهد.

\begin{matintable}{وضعیت قابلیت‌های کتاب‌نامه}{tab:bibliography-status}
\begin{tblr}{
    width = \textwidth,
    colspec = {l X[c] X[c]},
    hlines, vlines,
    colsep = 8pt,
    rowsep = 4pt,
    row{1} = {font=\bfseries},
}
قابلیت & توضیحات & وضعیت \\
\lr{\textbackslash cite} & استناد ساده به منابع & \textcolor{matingreen}{فعال} \\
\lr{\textbackslash parencite} & استناد پرانتزی & \textcolor{matingreen}{فعال} \\
\lr{\textbackslash textcite} & استناد متنی & \textcolor{matingreen}{فعال} \\
استناد چندگانه & چند منبع در یک استناد & \textcolor{matingreen}{فعال} \\
\lr{\textbackslash printbibliography} & چاپ کتاب‌نامه & \textcolor{matingreen}{فعال} \\
شیوه‌ی عددی & \lr{numeric} با ترتیب استناد & \textcolor{matingreen}{فعال} \\
عنوان فارسی & «منابع و مراجع» & \textcolor{matingreen}{فعال} \\
فایل \lr{.bib} & مدیریت منابع خارجی & \textcolor{matingreen}{فعال} \\
\end{tblr}
\end{matintable}

\section{معیارهای ارزیابی}
\label{sec:summary-criteria}

در این آزمون، سیستم کتاب‌نامه‌ی \lr{MatinBook} بر
اساس پنج معیار ارزیابی شد:

\begin{enumerate}
    \item \textbf{دستورات استناد}: چهار دستور اصلی.
    \item \textbf{استناد چندگانه}: \lr{\{a,b\}}.
    \item \textbf{نمایش کتاب‌نامه}: بدون \lr{latin}.
    \item \textbf{عنوان فارسی}: «منابع و مراجع».
    \item \textbf{چرخه‌ی biber}: کامپایل کامل.
\end{enumerate}

\section{نتیجه‌گیری}
\label{sec:summary-conclusion}

\textbf{تمام زیرسیستم‌های کتاب‌نامه‌ی \lr{MatinBook}
نسخه‌ی \lr{\matinbookversion} با موفقیت بارگذاری و
آزمایش شدند. تمام دستورات استناد و نمایش کتاب‌نامه
با \lr{biblatex} و \lr{biber} به‌درستی کار می‌کنند
و کتاب‌نامه با عنوان فارسی «منابع و مراجع» نمایش
داده می‌شود.}

%------------------------------------------------------------------------
% BACK MATTER (symmetric margins)
%------------------------------------------------------------------------
\backmatter
\frontmattergeometry

% Bibliography — NOT wrapped in \begin{latin}
\printbibliography[title={منابع و مراجع}]

% Index
\printindex

%------------------------------------------------------------------------
% BACK COVER
%------------------------------------------------------------------------
\authorbio{%
    تیم توسعه‌ی \lr{MatinBook}، گروهی از پژوهشگران و برنامه‌نویسان
    فارسی‌زبان است که با هدف ارائه‌ی یک راه‌حل حرفه‌ای برای
    حروف‌چینی کتاب‌های فنی فارسی فعالیت می‌کند. این پروژه حاصل
    سال‌ها تجربه در حوزه‌ی نشر دیجیتال و برنامه‌نویسی است.
}

\makebackcover{%
    این سند، دوازدهمین آزمون از مجموعه‌ی آزمون‌های یکپارچگی
    \lr{MatinBook} نسخه‌ی \lr{\matinbookversion} است.

    \vspace{0.5cm}

    \textbf{زیرسیستم‌های آزموده‌شده:}
    \begin{itemize}
        \item فایل \lr{.bib} خارجی
        \item دستورات استناد
        \item استنادهای چندگانه
        \item \lr{\textbackslash printbibliography}
        \item چرخه‌ی \lr{biber}
        \item ورودی‌های فارسی
    \end{itemize}
}

\end{document}